% 《大学物理学学习指导》自测题10填空10.2.1—10.2.8；
% 原书印刷207—208，PDF物理215—216页。
\subsection{自测题10：衍射与偏振填空}

\sourcetest{10.2.1}{207}{215}
\problemtext 惠更斯原理引入什么概念？
菲涅耳用什么思想补充惠更斯原理？
\workings
惠更斯把波前上的每一点看作产生\textbf{子波}的源。
菲涅耳再引入\textbf{子波相干叠加／干涉}的思想，
使波前次级波的叠加能解释衍射明暗。
\crosscheck 仅使用惠更斯几何包络而没有相干相位，
不能定量计算干涉衍射暗纹。

\sourcetest{10.2.2}{207}{215}
\problemtext 夫琅禾费单缝衍射，屏上P为第二级暗纹。
问对应波面可划成几个半波带；
若单缝宽变成一半，P将成为哪一级明/暗纹？
\workings
原来第二级暗纹满足
$a\sin\theta_P=2\lambda$，
所以全缝光程差含4个半波带。
缝宽改为$a/2$时同方向
\[
\frac a2\sin\theta_P=\lambda,
\]
对应第一级暗纹。故填：
\[
\boxed{N_{\frac12\lambda}=4},\qquad
\boxed{m_{\rm dark}=1}.
\]
\crosscheck 源题固定屏幕P与透镜光学几何，
半缝后相位差减半，并非把屏上坐标减半。

\sourcetest{10.2.3}{208}{216}
\problemtext 可见光400—760nm垂直入射光栅，
哪一级可见光谱可以在不与紧邻高级光谱重叠的条件下
保持完整？
\workings
第$m$级的红端对应$m(760)$；
第$m+1$级紫端对应$(m+1)(400)$。
若完整第$m$级无重叠，要求
\[
m(760)<(m+1)(400)
\ \Longleftrightarrow\ 360m<400.
\]
正整数$m$仅有$m=1$满足，故
\[
\boxed{\text{第一级光谱}}.
\]
\crosscheck 该判据针对相邻正级谱线在角位置的重叠，
还要求器件有足够的传播角范围才能观察其完整谱段。

\sourcetest{10.2.4}{208}{216}
\problemtext 多缝光栅常量$d=3\,\mu$m，透光缝宽$a=1\,\mu$m；
单缝衍射中央明纹内有多少条光栅主极大谱线？
\workings
光栅主极大$d\sin\theta=m\lambda$，
中央单缝包络界限$|a\sin\theta|<\lambda$。
故$|m|<d/a=3$，仅
$m=-2,-1,0,1,2$。
\[
\boxed{N=5}.
\]
\crosscheck $m=\pm3$正落单缝第一暗纹，不计入亮纹。

\sourcetest{10.2.5}{208}{216}
\problemtext 利用理想偏振片把一束入射线偏振光的振动方向旋转90°，
至少需要多少片偏振片？此时最大输出光强占输入的多少？
\workings
一片最终透振轴垂直入射振动方向的偏振片不能透光。
至少需一个中间轴和一个最终轴，即\textbf{两片}。
取中间角$\theta$，末片轴为90°，
\[
\frac{I_{\rm out}}{I_{\rm in}}=
\cos^2\theta\sin^2\theta
=\frac14\sin^22\theta\le\frac14.
\]
最大时$\theta=45^\circ$，
\[
\boxed{N_{\min}=2},\qquad\boxed{I_{\max}/I_{\rm in}=1/4}.
\]
\crosscheck 题目入射已经是线偏振光，
首片不需要额外乘自然光的$1/2$。

\sourcetest{10.2.6}{208}{216}
\problemtext 太阳光从空气射到平静水面$n=1.33$，
若反射光完全线偏振，求太阳光的仰角及反射光电矢量方向。
\workings
布儒斯特入射角对法线测量：
$i_B=\arctan n$。
太阳仰角是对水平水面测量，故
\[
\boxed{\alpha=90^\circ-\arctan(1.33)
\simeq36.94^\circ}.
\]
反射光仅有$s$偏振，
电场\textbf{垂直于入射面}。
源图中的入射面位于纸面内，故E矢量
\textbf{垂直纸面}。
\crosscheck 入射角与太阳仰角互余，不可将53.06°直接填作仰角。

\sourcetest{10.2.7}{208}{216}
\problemtext 图10.2.2展示五种光入射不同折射率界面的情况，
其中四幅为线偏振入射、一幅为自然光；
需在原图补出反射光与折射光，标注电矢量方向。
\workings
设入射面为图示纸面，
用$s$表示电场垂直入射面，
$p$表示电场平行入射面。
各向同性非磁界面的菲涅耳方程不会把两分量相互转换。
\[
\begin{array}{c|c|c}
\text{入射}&\text{反射}&\text{透射}\\ \hline
s\text{线偏振}&s\text{线偏振}&s\text{线偏振}\\
p\text{线偏振},\ i\ne i_B&p\text{线偏振}&p\text{线偏振}\\
p\text{线偏振},\ i=i_B&\text{反射消失}&p\text{线偏振}\\
\text{自然光},\ i\ne i_B&\text{部分偏振}&\text{部分偏振}\\
\text{自然光},\ i=i_B&s\text{线偏振}&\text{部分偏振}
\end{array}
\]
其中
$\boxed{i_B=\arctan(n_2/n_1)}$。
\crosscheck \textbf{原图逐线矢量重绘待二校：}
此处已给出决定源图A—E各图的物理分类准则，
但未将五个原版权子图重新描绘成可编辑原创TikZ示意，
因此本题的\emph{图形填绘部分未完成最终图解质检}。
不能把“原则正确”伪称每条箭头都已逐图校验。

\sourcetest{10.2.8}{208}{216}
\problemtext 在光学各向异性晶体内，
沿某特殊方向正常光与非常光有何相位传播特征？
该方向叫什么？只存在一个这种方向的晶体叫什么？
\workings
沿\textbf{光轴}方向，常光与非常光
的相位折射率相同、传播方向不再发生双折射分离。
只有一个光轴方向的晶体称为
\[
\boxed{\text{单轴晶体}}.
\]
\crosscheck 光轴是晶体内\emph{特殊传播方向}，
不是指穿过晶体几何中心的一条固定直线；
非晶各向同性介质没有相同的双折射结构。
